\scanpage{198}

haut poids pouvant figurer dans $\xi^n$ est $n$ fois le plus haut poids figurant dans $\xi$, soit $pn$. On a donc $pn+\alpha\leq pn$, d'où $\alpha\leq0$.

D'autre part, les poids intervenant dans $\xi$ sont $0$ et $p$, donc positifs, donc ceux dans $\xi^n$ sont positifs ; il faut donc qu'il \deleted{\unsure{$\alpha$}} en \deleted{soit} de $=$ des poids intervenant dans $\sigma_\alpha e_{0,p+1}$, d'où $\alpha\geq0$. D'où $\alpha=0$, i.e. $\sigma_\alpha=1$, c.q.f.d.

Ainsi :

% vertical rule at the left of the proposition
\noindent\begin{tabular}{@{}|p{.92\linewidth}@{}}\uline{Proposition} $A(G)$ est l'algèbre de polynômes sur $\Lambda=Z(\Gamma)$ $(\Gamma=Z)$ engendrée par $\xi=e_{0\,2}$ ; l'augmentation de $\Lambda$ est son augmentation usuelle, et $d(\xi)=\deleted{p}2$. Sur $\Lambda$, l'involution $\vee$ est l'involution usuelle, et on a
\[ \check\xi=e_{-p,2}=\sigma_{-p}.\xi. \]
\end{tabular}

% archive sheet number: 197
